\section{它为何为真？跳动的心脏}\label{sec:3}

我们已经看到了 $\Bun_{G}$ 的大致形状。但它究竟特殊在哪里？为什么要在它上面做范畴化调和分析？关键的想法是：

\begin{quote}
\itshape
\textbf{几何朗兰兹的精髓：}线性化之后，$\Bun_{G}$ 的行为惊人地像一个阿贝尔群。
\end{quote}

关键在于：$\Bun_{G}$ 那些神奇的性质并不出现在考察其几何、即在点的层面时，而出现在它的调和分析或量子力学、即线性化之后，具体说，是在过渡到层的范畴之后。这个口号有两个深刻的体现，它们指引着整个领域：

\begin{itemize}
\item \textbf{Hecke 算子}：$\Bun_{G}$ 上的层带有一大批彼此交换的对称性，以卷积的方式作用。
\item \textbf{Hitchin 纤维化}：在微局部化、即从位置空间过渡到相空间之后，$\Bun_{G}$ 在一般情形下看起来像雅可比簇。
\end{itemize}

我们先讲第一个，它在根本层面上解释了对偶群与对偶本身从何而来。第二个我们留到 \S\ref{sec:4}，它要几何得多（代价是只在一般情形下有效），让我们能具体看到本征层长什么样。

\subsection{局部对称性：Hecke 算子}\label{sec:3.1}

要理解 $\Bun_{G}$ 的对称性，回忆 $\mathrm{Pic}(C)$ 上群结构的局部起源——除子相加——是有益的。给定线丛 $L$ 与点 $x\in C$，我们可以只在 $x$ 处修改 $L$：考虑被要求以某阶消失（或被允许以某阶取极点）的截面。换言之，我们可以给表示 $L$ 的除子 $\sum n_{i}x_{i}$ 加上 $x$ 的一个倍数。这给出 $\Z$ 的一个拷贝通过 $L\mapsto L(n\cdot x)$ 作用在 $\mathrm{Pic}(C)$ 上，而这些作用（对变化的 $x$）生成整个群结构。

这个 $\Z$ 的非阿贝尔对应物是什么？给定 $G$-丛 $P\in\Bun_{G}$ 与点 $x\in C$，我们首先应当参数化 $P$ 在 $x$ 处的所有修改，即在 $x$ 之外与 $P$ 等同的丛。这样的修改由双陪集 $\Gr=\LGp\backslash\LG/\LGp$ 刻画，即仿射 Grassmannian $\Gr$ 上的 $\LGp$-轨道。高斯消元的变体表明，这些轨道的集合与对偶环面 $\check T$（$G$ 的环面 $T$ 的对偶）的特征标在 Weyl 群作用下的轨道一一对应（对 $G=\mathrm{GL}_{1}$ 就还原为我们的 $\Z$）。换言之，由最高权理论，这些轨道与 Langlands 对偶群 $\Gd$ 的不可约复表示一一对应；$\Gd$ 正是按组合方式定义的、以 $\check T$ 为对偶环面且与 $G$ 有相同 Weyl 群的群。

但这些丛的修改不只是一个集合：修改可以复合。确实，群中的双陪集 $K\backslash G/K$ 本身并不带结合乘法，但线性化之后就带了。也就是说，群上的乘法通过卷积在群元素的加权组合（即群上的测度）上定义了一个结合代数结构（群代数）。这个结构在过渡到双陪集时仍然保留：群上的卷积在 $K$-双不变测度上给出一个结合乘法，即 $G$ 与 $K$ 的 Hecke 代数。

我们关心的是用层（而不是函数）来线性化群，即用 D-模或可构造层。把这一想法用到双陪集空间 $\Gr$ 上，就得到了球 Hecke 范畴
\[
\Sph_{G}=\Shv(\Gr)=\Shv(\Gr)^{\LGp},
\]
即 $\LG$ 上关于 $\LGp$ 的双陪集上的所有层构成的范畴；更具体（但不那么对称地），它是仿射 Grassmannian $\Gr$ 上关于 $\LGp$ 作用等变的层，或更具体地，$\Gr$ 上在每个 $\LGp$-轨道上为常值的层。换言之，Hecke 范畴中的对象是丛在一点处的各种可能修改的加权和，权重是向量空间。Hecke 范畴 $\Sph_{G}$ 从修改的复合（最终来自 $\LG$ 上的卷积）继承一个结合乘法——它构成一个幺半范畴。此外，$\Sph_{G}$ 通过卷积算子，即 Hecke 函子，作用在 $\Shv(\Bun_{G})$ 上：我们把层在 $P\in\Bun_{G}$ 处的值替换为它在该点处由修改得到的诸丛上的加权平均。这些是图 Laplacian 的近亲——把函数在图上的值替换为它在邻居上的平均。

在~[KW07] 的规范场论描述中——那是在量子场论的线性化框架中进行的——Hecke 函子被实现为 't Hooft 线算子：$P$ 在 $x$ 处的修改是通过在 $x$ 处的规范场中制造一个指定的奇异性来实现的，这对应于制造一个穿过 $x$ 的磁单极子。

\subsection{因子化：交换性之源}\label{sec:3.2}

GLC（可以说也是整个 Langlands 纲领）确凿无疑的跳动心脏是 Hecke 修改的奇妙交换性。一般的 Hecke 代数作为非阿贝尔群乘法的平均版本，本来没有理由交换。但磁单极算子的物理图像提示了另一种复合图景，即算子乘积展开（OPE）：我们可以在不同的、但彼此邻近的点上作相继的观测，然后取两点碰撞的极限。与量子力学本质上的非交换性不同，在更高维的量子场论中，算子有更多空间可以彼此绕行。在拓扑量子场论（测量沿形变保持不变）的框架中，这使算子的乘积交换。

同样的想法在代数拓扑中很熟悉：拓扑空间的高阶同伦群 $\pi_{2}$（或对 $n\ge2$ 的 $\pi_{n}$）之所以交换、而 $\pi_{1}$ 不交换，正是因为在一个带基点的空间中，我们可以让两个二维环路——把圆盘边界送到基点的映射——在更大圆盘的映射内部彼此绕行，这使得它们的复合在同伦意义下交换。这一观察是小圆盘（或 $E_{n}$）operad（乃至一般 operad）理论的起源，也是同伦论中交换性的主要来源。

这与我们何干？仿射 Grassmannian $\Gr$ 事实上同伦等价于分类空间 $BG$ 的二重环路空间：它被定义为曲面上的 $G$-丛空间，细看之下它描述的是在小圆盘外平凡化的丛，因而是到 $BG$ 的带基点映射。这意味着 $\Gr$ 本身带有同伦交换的乘法。

Beilinson 与 Drinfeld~[BD] 发现了这些来自物理与同伦论的机制的一个几何统一，称为因子化（factorization，同义地也称融合 fusion），它在 Hecke 范畴 $\Sph_{G}$ 背后起作用。也就是说，我们不必在同一点 $x\in C$ 处对丛作相继的 Hecke 修改（即用两个 Hecke 算子 $T_{1},T_{2}$ 作用），而可以在 $x$ 处作修改 $T_{1}$、在邻近的 $y$ 处作 $T_{2}$。球范畴的拓扑本性保证：当我们移动 $y$、或取 $y$ 与 $x$ 碰撞的极限时，复合在局部是常值的。让 $y$ 完全绕过 $x$ 一周，我们就得到了从复合 $T_{1}\circ T_{2}$ 到 $T_{2}\circ T_{1}$ 的一个形变。结论是：$\Sph$ 上的卷积乘积实际上是交换的——$\Sph$ 构成一个对称幺半范畴。

因子化——一种纯粹由流形上点的构型及其碰撞来几何地参数化、而非由传统的复合映射给出的代数结构——是一次真正革命性的发现，是几何朗兰兹研究中产生的最深刻、影响最大的思想。它的源头是 Segal、Beilinson 与 Drinfeld 为理解二维共形场论中算子乘积展开的几何意义所作的努力（这些在代数上被编纂为顶点代数与辫张量范畴的理论，那里“融合”一词更常用），以及它与代数拓扑中迭代环路空间（如 $\Gr$）与小圆盘 operad 的代数理论之间浮现出的密切平行。

在~[BD04] 中编纂之后，它为物理、拓扑、范畴论与数论中的现象提供了一个统一的视角。反过来说，在 Costello 与 Gwilliam~[CG17] 发展的一般数学框架中，因子化正是量子场论中可观测量的定义性特征。

Hecke 算子的交换性可以说是 Langlands 纲领最不可或缺的一面——没有可对角化的交换算子，就没有自守形式的谱分解。它背后是自守 $L$-函数的构造，并提供了把表示论与 Galois 理论联系起来的根本机制。然而这一交换性是通过一个技巧证明的，这个技巧既不能解释它为何成立，也不能说明如何把它变成一个稳固而可塑的原理。因子化恰好提供了这个原理，由它可以推出经典的交换性以及更多东西（通过去范畴化）。得益于 V.~Lafforgue、Fargues、Scholze、Zhu 以及许多其他人的开创性工作，它已成为算术背景中的基本工具。

\subsection{对偶群的诞生}\label{sec:3.3}

既然知道球 Hecke 范畴是交换的，我们就应当寻找 $\Sph$ 的傅里叶对偶描述——用张量积而非卷积的语言、即使其交换性显而易见的描述。更几何地说，我们应当寻找它的谱，即其作用的所有可能（联合）本征值构成的空间。这正是几何 Satake 对应的内容。已经有一个自然的猜测：我们注意到双陪集 $\Gr=\LGp\backslash\LG/\LGp$ 与 Langlands 对偶约化群 $\Gd$ 的不可约表示一一对应。于是人们可以期望用 $\Gd$-表示及其张量积来描述 $\Sph_{G}$ 的整个结构。下面的定理证实了这一猜测：

\begin{theorem*}[几何 Satake 对应；[Lus83, Gin, MV07]]\label{thm:satake}
存在对称幺半阿贝尔范畴的等价
\[
\Sph_{G}\;\simeq\;\Rep(\Gd),
\]
它把带卷积的球范畴与带张量积的 Langlands 对偶群 $\Gd$ 的表示范畴等同起来。
\end{theorem*}

\subsection{注}\label{sec:3.4}

该定理在~[MV07] 中被推广到系数取自任意域、甚至整数环的（反）周层（以及 $\Gd$ 的表示）。为简单起见，我们隐含地从导出范畴过渡到它们的核心（heart），即阿贝尔 $\C$-线性范畴，于是 $\Sph^{\heartsuit}$ 等价地由仿射 Grassmannian 上 $\LGp$-等变的 D-模或（反）周层构成。几何 Satake 还有对完整导出范畴的、重要而微妙得多的版本~[BF08]。

几何 Satake 对应或许最好反过来理解：它提供了 Langlands 对偶群 $\Gd$ 的真正定义。正如交换环定义了一个几何对象——它的谱，参数化该交换环所有可能的联合本征值——像 $\Sph$ 这样的交换张量范畴也有一个由 Tannaka 重构给出的谱。这参数化了所有可能的联合本征值，只不过这些本征值现在是向量空间而非算子。确实，对 $G=\Gd=\mathrm{GL}_{n}$，该定理可以转述为一个等价
\[
\{\Sph_{\mathrm{GL}_{n}}\ \text{的可能本征值}\}\;\longleftrightarrow\;\{n\ \text{维向量空间}\}.
\]
$n$ 维向量空间在同构意义下只有一个，但定理抓住了它的自同构群，即 $\Gd=\mathrm{GL}_{n}$。换言之，右端代表 $n$ 维向量空间的模栈 $\mathrm{pt}/\mathrm{GL}_{n}$。另一方面，对 $G=\mathrm{SO}_{2n+1}$（奇正交群），对偶群为 $\Gd=\mathrm{Sp}_{2n}$，定理说的是
\[
\{\Sph_{\mathrm{SO}_{2n+1}}\ \text{的可能本征值}\}\;\longleftrightarrow\;\{2n\ \text{维辛向量空间}\}.
\]
对一般的 $G$，答案由一点上的 $\Gd$-局部系统给出——那正是把 $\Gd$ 的任意表示变成向量空间、从而把 Hecke 算子变成本征值的机器。\zhnote{几何 Satake 对应的内容与 nLab 词条一致：$\Sph_{G}$ 是 $\LGp$-等变层构成的交换张量范畴，其 Tannaka 对偶（谱）正是 $\Rep(\Gd)$，于是对偶群 $\Gd$ 由 $\Sph_{G}$ 的自同构群重建出来；这正是 Tannaka 重构的几何版本，见 \emph{geometric Satake equivalence} 与 \emph{Tannaka duality} 两个词条：\href{https://ncatlab.org/nlab/show/geometric+Satake+equivalence}{\texttt{https:/\allowbreak{}/\allowbreak{}ncatlab.\allowbreak{}org/\allowbreak{}nlab/\allowbreak{}show/\allowbreak{}geometric+\allowbreak{}Satake+\allowbreak{}equivalence}}、\href{https://ncatlab.org/nlab/show/Tannaka+duality}{\texttt{https:/\allowbreak{}/\allowbreak{}ncatlab.\allowbreak{}org/\allowbreak{}nlab/\allowbreak{}show/\allowbreak{}Tannaka+\allowbreak{}duality}}（最后修订 2025-07-28）。}

为证明几何 Satake，我们先构造一个群，作为 $\Sph_{G}$ 上总上同调函子的自同构群（这个函子最终将对应于遗忘函子 $\Rep(\Gd)\to\mathrm{Vect}$）。然后有一个从 $\Sph$ 到该群表示的平凡张量函子，我们验证它是等价。这严重依赖于 $\Sph_{G}$ 的交换性，从而依赖于因子化。Grassmannian 及其 Schubert 分层的优美几何随后迫使这个群是约化的，并让我们算出其根系数据在组合意义下与 $G$ 的对偶，于是我们就有权把它叫作 Langlands 对偶群。

\subsection{整体对称性：谱作用}\label{sec:3.5}

借助几何 Satake 对应，我们现在有一大批交换算子作用在 $\Shv(\Bun_{G})$ 上——对每个点 $x\in C$，我们都得到一个交换张量范畴 $\Rep(\Gd)$ 的拷贝，而且它们彼此交换。换言之，我们找到了除子群 $\mathrm{Div}(C)=\bigoplus_{x\in C}\Z$ 作用在 Picard 群上的非阿贝尔对应物（这里的 $\Z$ 现在被理解为标记 $\mathrm{GL}_{1}(\C)\cong\C^{\times}$ 的不可约表示）。这提示我们也应找到除子之间作用于 Picard 群上的关系的对应物，即不同点处的修改如何相互作用。从谱的角度看，几何 Satake 描述的是单点处 Hecke 算子的谱，而现在我们需要理解由 $C$ 上所有点处的丛修改所生成的“整体 Hecke 代数”的谱。

关键仍然在因子化结构中：我们已经用一种把 $C$ 上邻近点的移动与碰撞包含在内的方式理解了 Hecke 算子的复合，所以只需理解从局部到整体的过渡。换言之，这一大批 Hecke 算子满足这样的关系：作用以局部常值的方式依赖于修改点 $x$，而在两个碰撞的点处作修改等同于在一点处复合修改。

这一非形式的想法被因子化代数的因子化同调这一概念精确捕捉。Beilinson–Drinfeld~[BD04] 的这一想法是顶点代数共形块理论的抽象化，如今为量子场论中整体可观测量如何由局部可观测量装配而成提供了一个有影响的数学模型~[CG17]，并在拓扑与高阶代数中作为经典同调理论的推广成为一个强大的多用途工具~[AF15, Lur]。\zhnote{因子化同调（factorization homology）是把因子化代数在流形上“积分”（即取张量积式的同调）所得的不变量；它在代数拓扑中的原型是 Ayala–Françis 对拓扑流形的构造。nLab 词条 \emph{factorization homology} 给出定义：对流形 $M$ 与因子化代数 $A$，$\int_{M}A$ 由把 $M$ 分解为小圆盘、在每块上取 $A(D)$、再沿嵌入取（高阶）余极限得到：\href{https://ncatlab.org/nlab/show/factorization+homology}{\texttt{https:/\allowbreak{}/\allowbreak{}ncatlab.\allowbreak{}org/\allowbreak{}nlab/\allowbreak{}show/\allowbreak{}factorization+\allowbreak{}homology}}（最后修订 2026-05-30）。原文强调的正是“局部数据 $\Rightarrow$ 整体可观测量”的装配机制。}

于是现在的挑战，与几何 Satake 对应时一样，是以对角化或傅里叶对偶的方式、用张量积而非卷积来描述球 Hecke 范畴的整体化（因子化同调）。换言之，我们要描述整体 Hecke 范畴的谱，即它所有可能的本征值构成的汇集。

自然的猜测又出现了。先考虑 $G=\mathrm{GL}_{n}$ 的情形。如前所述，$\Sph_{\mathrm{GL}_{n}}$ 在单点 $x\in C$ 处作用的可能本征值由 $n$ 维向量空间给出。上面的讨论于是提示：整体上应当有一族向量空间 $\{E_{x}\}_{x\in C}$，它们在 $x$ 上局部常值，即邻近点处的向量空间相互等同。这正是一个秩 $n$ 的 $C$ 上局部系统的数据，即特征簇 $\Loc_{\mathrm{GL}_{n}}(C)$ 的一个点——不同构的局部系统由它们沿 $C$ 中环路的单值性区分。更一般地，$C$ 上的 $\Gd$-局部系统 $E$ 正是把 $\Gd$ 的表示 $V$ 变成 $C$ 上局部常值向量空间族 $(V)_{E}$ 的机器。

这引向下述基本定理：

\begin{theorem*}[谱作用]\label{thm:spectralaction}
Hecke 作用 $\bigotimes_{x\in C}\Sph_{x}$ 在 $\Shv(\Bun_{G})$ 上的作用通过 $\QCoh(\Loc_{\Gd}(C))$ 的作用分解\anote{确切地说，映射 $\bigotimes_{x\in C}\Rep(\Gd)\to\QCoh(\Loc_{\Gd})$ 把 $x\in C$ 处的表示 $V\in\Rep(\Gd)$ 送到 $\Loc_{\Gd}$ 上的向量丛，其在局部系统 $E$ 处的纤维为 $(V)_{E}|_{x}$。}，de Rham~[Gai15] 与 Betti~[NY19] 两种背景均成立。
\end{theorem*}

与几何 Satake 对应一样，我们可以把这个结果转述为
\[
\{\Shv(\Bun_{G})\ \text{上 Hecke 算子的可能联合本征值}\}\;\longleftrightarrow\;\{C\ \text{上的}\ \Gd\text{-局部系统}\}.
\]

谱作用提供了几何朗兰兹对应中自守到谱的方向——它告诉我们，自守范畴可以在 Langlands 参数空间上作谱分解。遵循同样总体思路的谱作用变体现已在许多不同背景下得到证明，并构成本领域的若干最伟大成就，包括函数域上的 Langlands 对应~[Laf18, AGK+20] 与非阿基米德局部域上的局部 Langlands 对应~[FS21]。\zhnote{谱作用定理的“谱侧”是 $\QCoh(\Loc_{\Gd})$，其中 $\Loc_{\Gd}$ 是 $C$ 上 $\Gd$-局部系统的模栈；“作用”来自每个点 $x$ 处的几何 Satake 等价与 $\Rep(\Gd)$ 的对称幺半结构，经由因子化同调粘成整体。nLab 词条 \emph{geometric Langlands correspondence} 在“谱作用”一节给出的陈述与此一致：存在 $\QCoh(\Loc_{\Gd})$ 在 $\Shv(\Bun_{G})$ 上的作用，其单点情形即几何 Satake 等价：\href{https://ncatlab.org/nlab/show/geometric+Langlands+correspondence}{\texttt{https:/\allowbreak{}/\allowbreak{}ncatlab.\allowbreak{}org/\allowbreak{}nlab/\allowbreak{}show/\allowbreak{}geometric+\allowbreak{}Langlands+\allowbreak{}correspondence}}（最后修订 2026-04-11）。}

从规范理论的观点看，几何 Satake 对应把 $G$-丛理论中的自然局部可观测量（Hecke 磁单极算子）与对偶 $\Gd$-规范理论中由 $\Gd$-局部系统的单值性构造出的那些（Wilson 环路算子）等同起来。谱作用于是断言 $C$ 上的整体可观测量作用在 $\Bun_{G}$ 上的自守层上——换言之，它在可观测量的层面建立了电磁对偶的一个版本。
